\documentclass[a4paper, 10pt]{article}
\usepackage{scrextend}
\changefontsizes{9pt}
% Per-subject config for this repo. See vendor/template/course-config.tex.example
% for what each macro is used for.

\newcommand{\CourseCode}{2110506}
\newcommand{\CourseName}{Software-Defined Systems}
\newcommand{\StudentID}{6733172621}
\newcommand{\StudentName}{Patthadon Phengpinij}
\newcommand{\DefaultCollaborators}{Claude (for\,\LaTeX\,styling and grammar checking)}
\newcommand{\AssignmentLabel}{Activity}

\usepackage{../vendor/template/CEDT-Assignment-style}

\begin{document}
\subject
\asgntitle{1}{}{Week 2}{\StudentID\ \StudentName}{\DefaultCollaborators}

% ================================================================================ %
\section*{Activity: Docker}
% ================================================================================ %

\textthai{จากตัวอย่างในชั้นเรียนที่ทำ} multi-container application \textthai{ซึ่งประกอบด้วย} container node-exporter,
Prometheus \textthai{และ} Grafana \textthai{ให้นิสิตเพิ่ม} container
\begin{enumerate}
	\item Apache web server
	\item apache-exporter \textthai{ซึ่งทำหน้าที่เก็บข้อมูลจาก} Apache web server
\end{enumerate}

\par\noindent \underline{\textthai{การตั้งค่า} Apache Web server \textthai{เพื่อให้สร้างหน้า} server status}
\par \textthai{ให้สร้างไฟล์ชื่อ} status.conf
\begin{codingbox}[language=bash]
<IfModule mod_status.c>
# Allow server status reports generated by mod_status,
# with the URL of http://servername/server-status
# Uncomment and change the "192.0.2.0/24" to allow access from other hosts.

<Location /server-status>
SetHandler server-status
<RequireAll>
Require all granted
</RequireAll>
</Location>

# Keep track of extended status information for each request
ExtendedStatus On

# Determine if mod_status displays the first 63 characters of a request or
# the last 63, assuming the request itself is greater than 63 chars.
# Default: Off
#SeeRequestTail On

<IfModule mod_proxy.c>
# Show Proxy LoadBalancer status in mod_status
ProxyStatus On
</IfModule>

</IfModule>
\end{codingbox}

\par\noindent \textthai{สร้าง} network \textthai{เพื่อให้} container \textthai{ทุกตัวติดต่อกันด้วยชื่อ} (\textthai{หรือใช้} mynet \textthai{จากตัวอย่างในชั้นเรียน}):

\begin{center}\texttt{docker network create monitoring}\end{center}

\par\noindent \textthai{\textbf{อย่าลืม}เชื่อม} container (node-exporter, Prometheus, Grafana)
\textthai{เข้ากับ} network \textthai{นี้ด้วย} (\texttt{--network monitoring} \textthai{หรือ} \texttt{docker network connect})

\par\noindent \textthai{สร้าง} container \textthai{โดยใช้คำสั่งดังนี้}:
\begin{codingbox}[language=bash]
$ docker run -d --name apache --network monitoring -p 8080:80 -v "$(pwd)/status.conf:/etc/apache2/modsenabled/status.conf" ubuntu/apache2:latest
\end{codingbox}

\par\noindent \textbf{\textthai{หมายเหตุ}}: \textthai{ต้องมีไฟล์} \texttt{status.conf}
\textthai{อยู่ใน} directory \textthai{ปัจจุบันก่อนรัน} (\textthai{ใช้} absolute path)

\par\noindent \textthai{เมื่อเปิดหน้าเว็บที่} \url{http://localhost:8080/server-status/?auto}
\textthai{จะได้ผลลัพธ์ดังตัวอย่าง}
\fig[0.5\textwidth]{images/server-status-example.png}{\textthai{ตัวอย่างผลลัพธ์ของหน้า} server status}{http://localhost:8080/server-status/?auto}

\par\noindent \textthai{การสร้าง} container \textthai{ของ} Apache exporter \textthai{ให้รันคำสั่งดังตัวอย่าง}:
\begin{codingbox}[language=bash]
$ docker run -d --name apache-exporter --network monitoring -p 9117:9117 lusotycoon/apache-exporter --scrape_uri="http://apache/server-status?auto"
$
\end{codingbox}

\par\noindent \textthai{เมื่อไปที่} \url{http://localhost:9117/metrics} \textthai{จะได้ผลลัพธ์ดังตัวอย่าง}
\fig[0.8\textwidth]{images/apache-exporter-metrics-example.png}{\textthai{ตัวอย่างผลลัพธ์ของหน้า} Apache exporter metrics}{http://localhost:9117/metrics}

\par\noindent \textthai{จากนั้นให้นิสิตแก้ไขไฟล์} \texttt{prometheus.yml}
\textthai{เพื่อเพิ่ม} apache \textthai{และ} apache exporter \textthai{เข้าไปในระบบ}
\textthai{โดยเพิ่ม} job \textthai{ต่อไปนี้ใน} section scrape\_configs
(\textthai{สังเกตว่า} Prometheus \textthai{ดึงข้อมูลจาก} exporter \textthai{ที่} port 9117
\textthai{ไม่ได้ดึงจาก} Apache \textthai{โดยตรง}):

\begin{codingbox}[language=yaml]
scrape_configs:
 # ... existing node-exporter job ...
 - job_name: "apache"
 static_configs:
 - targets: ["apache-exporter:9117"]
\end{codingbox}

\par\noindent Prometheus \textthai{สามารถเข้าถึง} apache-exporter
\textthai{ด้วยชื่อ} container \textthai{บน} network (\texttt{apache-exporter:9117})
\textthai{โดยในส่วนของคำสั่ง} \texttt{-p 9117:9117} \textthai{มีไว้ให้เปิดดูจาก}
browser \textthai{บน} host \textthai{เท่านั้น เมื่อแก้ไขแล้วให้} reload Prometheus
\textthai{และตรวจสอบที่หน้า} Targets \textthai{ว่ามีสถานะ} UP

\myspace

\par\noindent \underline{\textbf{\textthai{สิ่งที่ต้องส่งได้แก่}}}
\begin{enumerate}[label=\arabic*.]
	\item \textthai{ไฟล์} \texttt{prometheus.yml}
	\item \textthai{ภาพ} capture \textthai{หน้าจอ พร้อมคำอธิบายสั้น ๆ ประกอบแต่ละภาพของ}
	      \begin{enumerate}[label=\arabic*.]
		      \item \textthai{การรันคำสั่งและผลลัพธ์ที่ได้ของ}
		            \begin{enumerate}[label=\arabic*.]
			            \item \textthai{คำสั่งที่สร้าง} network
			            \item \textthai{คำสั่งที่สร้าง} persistent storage \textthai{ได้แก่} volume \textthai{ของ} Grafana \textthai{จากตัวอย่างในชั้นเรียน}
			                  (\textthai{เพื่อให้} dashboard \textthai{คงอยู่หลังสร้าง} container \textthai{ใหม่})
			                  \textthai{และ} bind mount \textthai{ของไฟล์} \texttt{status.conf}
			            \item \textthai{คำสั่งที่รัน} container \textthai{แต่ละตัว}
		            \end{enumerate}
		      \item \textthai{หน้าจอ} Grafana \textthai{ที่มีกราฟจาก} node-exporter
		            \textthai{และ} apache-exporter \textthai{พร้อมคำอธิบายสั้น ๆ เกี่ยวกับ} metrics \textthai{ที่เลือกใช้}
	      \end{enumerate}

\end{enumerate}

% ================================================================================ %
%                                    Problem 01                                    %
% ================================================================================ %
\begin{problem}
\textthai{ไฟล์} \texttt{prometheus.yml}
\end{problem}

\begin{solution}
	Here's the content of the \texttt{prometheus.yml} file:
	\begin{codingbox}[language=yaml]
global:
  scrape_interval: 15s

scrape_configs:
- job_name: "node"
  static_configs:
    - targets: ["node-exporter:9100"]
- job_name: "apache"
  static_configs:
    - targets: ["apache-exporter:9117"]
\end{codingbox}
\end{solution}
% ================================================================================ %

% ================================================================================ %
%                                    Problem 02                                    %
% ================================================================================ %
\begin{problem}
\textthai{ภาพ} capture \textthai{หน้าจอ พร้อมคำอธิบายสั้น ๆ ประกอบแต่ละภาพของการทำงานในแต่ละขั้นตอน}
\end{problem}

\begin{solution}
	\begin{enumerate}[label=2.\arabic*.]
		\item \textthai{การรันคำสั่งและผลลัพธ์ที่ได้ของ}
		      \begin{enumerate}[label*=\arabic*.]
			      \item \textthai{คำสั่งที่สร้าง} network
			            \begin{codingbox}
$ docker network create monitoring
$
\end{codingbox}
			            \fig[0.8\linewidth]{images/screenshot-2.1.1.png}{Screenshot of the results after running the command to create the network.}{network}
			      \item \textthai{คำสั่งที่สร้าง} persistent storage \textthai{ได้แก่} volume \textthai{ของ} Grafana \textthai{จากตัวอย่างในชั้นเรียน}
			            (\textthai{เพื่อให้} dashboard \textthai{คงอยู่หลังสร้าง} container \textthai{ใหม่})
			            \textthai{และ} bind mount \textthai{ของไฟล์} \texttt{status.conf}
			            \par\noindent \textbf{Persistent storage for Grafana.}
			            \begin{codingbox}
$ docker volume create grafana-vol
$
\end{codingbox}
			            \fig[0.8\linewidth]{images/screenshot-2.1.2-1.png}{Screenshot of the results after creating persistent storage.}{grafana persistent-storage}
			            \par\noindent \textbf{Bind mount for Apache status configuration.}
			            \begin{codingbox}
$ docker run --rm -d -q -p 8080:80 --name apache --network monitoring -v "$(pwd)/code/status.conf:/etc/apache2/mods-enabled/status.conf" ubuntu/apache2:latest
\end{codingbox}
			            \fig[\linewidth]{images/screenshot-2.1.2-2.png}{Screenshot of the results after creating the bind mount for Apache status configuration.}{apache bind-mount}
			      \item \textthai{คำสั่งที่รัน} container \textthai{แต่ละตัว}
			            \begin{codingbox}
$ # Node Exporter
$ docker run --rm -d -q -p 9100:9100 --name node-exporter --network monitoring prom/node-exporter
$ # Prometheus
$ docker run --rm -d -q -p 9090:9090 --name prometheus --network monitoring -v "$(pwd)/code/prometheus.yml:/etc/prometheus/prometheus.yml" prom/prometheus:v3.0.1
$ # Grafana
$ docker run --rm -d -q -p 3000:3000 --name grafana --network monitoring -v grafana-vol:/var/lib/grafana grafana/grafana:11.4.0
$ # Apache
$ docker run --rm -d -q -p 8080:80 --name apache --network monitoring -v "$(pwd)/code/status.conf:/etc/apache2/mods-enabled/status.conf" ubuntu/apache2:latest
$ # Apache Exporter
$ docker run --rm -d -q -p 9117:9117 --name apache-exporter --network monitoring lusotycoon/apache-exporter --scrape_uri="https://apache/service-status?auto"
\end{codingbox}
									\fig[0.8\linewidth]{images/screenshot-2.1.3.png}{Screenshot of the results after creating all the containers.}{all containers}
		      \end{enumerate}
		\item \textthai{หน้าจอ} Grafana \textthai{ที่มีกราฟจาก} node-exporter
		      \textthai{และ} apache-exporter \textthai{พร้อมคำอธิบายสั้น ๆ เกี่ยวกับ} metrics \textthai{ที่เลือกใช้}
					\fig[0.8\linewidth]{images/grafana-visualization.png}{Screenshot of the Grafana dashboard showing metrics from node-exporter and apache-exporter.}{grafana dashboard}
					\textbf{Selected metrics from node-exporter and apache-exporter:}
					\begin{itemize}
						\item \textbf{Apache CPU Load:}
									\textthai{เลือกใช้} Apache CPU Load \textthai{เพื่อดูปริมาณการใช้งาน} CPU
									\textthai{ที่เกิดจากการทำงานของ} Apache Web Server
									\textthai{โดยตรง ช่วยให้สามารถตรวจสอบได้ว่าเมื่อมีผู้ใช้งานหรือ} Request \textthai{เข้ามามากขึ้น}
									Apache \textthai{มีการใช้ทรัพยากร} CPU \textthai{สูงขึ้นหรือไม่}
									\textthai{หากค่า} CPU Load \textthai{สูงต่อเนื่อง อาจบ่งบอกว่าเซิร์ฟเวอร์กำลังรับภาระหนักและอาจส่งผลต่อประสิทธิภาพในการให้บริการเว็บไซต์}
						\item \textbf{Node Load 1M:}
									\textthai{เลือกใช้} Node Load 1M \textthai{เพื่อดูค่า} Load Average
									\textthai{ของระบบในช่วง 1 นาทีล่าสุด ซึ่งแสดงภาพรวมของภาระงานบนเครื่องเซิร์ฟเวอร์แบบใกล้เคียงเวลาจริง}
									\textthai{ช่วยตรวจสอบได้อย่างรวดเร็วว่าระบบกำลังมีงานหรือ} Process \textthai{กำลังรอใช้} CPU
									\textthai{มากผิดปกติหรือไม่ จึงเหมาะสำหรับการเฝ้าระวังการเปลี่ยนแปลงของภาระระบบในระยะสั้น}
					\end{itemize}
	\end{enumerate}
\end{solution}
% ================================================================================ %

% ================================================================================ %
%                                    Problem 03                                    %
% ================================================================================ %
\begin{problem}
Screenshot of the results after running the \texttt{GRADER}.
\end{problem}

\begin{solution}
	\fig[\linewidth]{images/grader.png}{Screenshot of the results after running the \texttt{GRADER}.}{grader}
\end{solution}
% ================================================================================ %

\end{document}
